\documentclass[10pt,a4paper]{amsart}
\usepackage[margin=19mm]{geometry}
\usepackage{amsmath,amssymb,mathtools}
\usepackage[scr=rsfs,cal=euler]{mathalfa}
\usepackage{xfrac}
\usepackage[unicode,hidelinks,bookmarks=false]{hyperref}
\newcommand{\ie}{i.e.\ }
\newcommand{\cf}{cf.\ }
\newcommand{\resp}{respectively\ }
\newcommand{\Subsection}{\subsection}
\providecommand{\nobreakdash}{\nobreak\hbox{-}}
\setlength{\emergencystretch}{2em}
\multlinegap0pt
\usepackage{amssymb}
\usepackage{multirow}
\usepackage{bbm}
\usepackage{tikz}
\usepackage{float}
\usepackage{enumerate}
\usepackage{mathtools}
\newtheorem{lemma}{Lemma}
\newtheorem{theorem}{Theorem}
\title{Arithmetic Progressions of Integers that are\\Relatively Prime to their Digital Sums}
\author{Ryan Blau, Joshua Harrington, Sarah Lohrey, Eliel Sosis, Tony W. H. Wong}
\date{Source: arXiv:2308.05228v1 (CC BY 4.0)}
\begin{document}
\maketitle

\noindent\emph{Source and licence.} Arithmetic Progressions of Integers that are\\Relatively Prime to their Digital Sums. Version arXiv:2308.05228v1 (CC BY 4.0), distributed by arXiv under the Creative Commons Attribution 4.0 International licence, http://creativecommons.org/licenses/by/4.0/, as recorded in the arXiv licence metadata for this exact version and shown on the abstract page https://arxiv.org/abs/2308.05228v1. Full licence notice: \texttt{licenses/arxiv-2308.05228-CC-BY-4.0.txt}.\par\smallskip

\noindent\textit{Editorial scope.} Counted unit: the article's theorem thm:bevenb-1AP (source0.tex lines 263--265). The article's complete original proof of it is delivered with it (source0.tex lines 266--310, one \texttt{proof} environment of 45 lines). No mathematics is written, repaired or re-derived: the statement and the proof are the article's own lines, in the article's own notation, and no retained line is substituted or re-flowed.  Retained uncounted as the source-local context the proof needs: lines 219--221.  Nothing was omitted from any retained span, so the package carries no omission record.  No retained line contains a citation, so the wrapper carries no bibliography and no bibliography entry.  No retained line contains a figure environment, an image-inclusion command or drawing code, so no image is delivered and the package declares no asset dependency.  Editorial formatting only, in addition: an AMS \texttt{amsart} preamble that restates the article's own declarations and macros used by the retained lines, namely the environments \texttt{lemma}, \texttt{theorem} on separate counters, together with the standard packages those lines need.  The retained environments print in this document as Lemma 1, Theorem 1 in order of appearance, whereas the article prints the same items with the numbering of its own full document.  The complete native source of the article as distributed in the arXiv e-print for this exact version is bundled unchanged as \texttt{sources/145441/source0.tex}.
\begin{lemma}\label{lem:b^m=b}
For all finite collections of distinct primes $q_1,q_2,\dotsc,q_t$, there exist infinitely many positive integers $m$ such that $b^m\equiv b\pmod{q}$.
\end{lemma}

\begin{theorem}\label{thm:bevenb-1AP}
Let $b$ be even. Then the maximum length of a $b$-anti-Niven $(b-1)$-AP is $2b+1$. Furthermore, there exist infinitely many such sequences of length $2b+1$.
\end{theorem}

\begin{proof}
Suppose there is a $b$-anti-Niven $(b-1)$-AP $\mathcal{S}$ of length at least $2b+1$. Since $\gcd(b-1,b)=1$, there exist two terms in $\mathcal{S}$ that are multiples of $b$. Let these two terms be $ab$ and $ab+(b-1)b$ for some positive integer $a$. Note that $ab$ is even, so $s_b(ab)=s_b(a)$ is odd. Since $ab+2(b-1)=(a+1)b+b-2$ is an even term in $\mathcal{S}$, the digit sum $s_b((a+1)b+b-2)=s_b(a+1)+b-2$ must be odd, implying that $s_b(a+1)$ is also odd. Hence, $a=cb+b-1$ for some nonnegative integer $c$, where $s_b(c)$ is even. In other words, $ab=cb^2+(b-1)b$ and $ab+(b-1)b=(c+1)b^2+(b-2)b$. Also, $s_b(c+1)$ is odd since $s_b((c+1)b^2+(b-2)b)=s_b(c+1)+b-2$ is odd.

Now, note that $cb^2$ and $s_b(cb^2)=s_b(c)$ are even, so $cb^2$ is not in $\mathcal{S}$. Similarly, $(c+1)b^2+(b-1)b+b-2$ and $s_b((c+1)b^2+(b-1)b+b-2)=s_b(c+1)+b-1+b-2$ are even, so $(c+1)b^2+(b-1)b+b-2=cb^2+(2b+2)(b-1)$ is also not in $\mathcal{S}$. Therefore, $\mathcal{S}$ is a subsequence of $\{cb^2+j(b-1):1\leq j\leq2b+1\}$, thus the maximum length of a $b$-anti-Niven $(b-1)$-AP is at most is $2b+1$.

It remains to show that such sequences occur infinitely often. Let $c$ be a nonnegative integer such that $s_b(c+1)=s_b(c)+1$. Then it is not difficult to observe that $s_b(cb^2+j(b-1))=s_b(c)+b-1$ for $1\leq j\leq b$ and $b+2\leq j\leq 2b$, and $s_b(cb^2+j(b-1))=s_b(c)+2(b-1)$ for $j\in\{b+1,2b+1\}$. Hence, it suffices to show that there exist infinitely many positive integers $c$ such that
\begin{itemize}
%\item $s_b(c)$ is even,
\item $b\nmid(c+1)$,
\item $\gcd(cb^2+j(b-1),s_b(c)+b-1)=1$ for $1\leq j\leq b$ and $b+2\leq j\leq 2b$, and
\item $\gcd(cb^2+j(b-1),s_b(c)+2(b-1))=1$ for $j\in\{b+1,2b+1\}$.
\end{itemize}
Let $p_1,p_2,\dotsc,p_t$ be all primes less than or equal to $2b$. By Lemma~\ref{lem:b^m=b}, there exist infinitely many positive integers $m$ such that $b^{m+1}\equiv b\pmod{p_1p_2\dotsb p_t}$. In other words, $p_1p_2\dotsb p_t\mid b(b^m-1)$. Let $P=b^m+1$. Since $\gcd(b^m+1,b)=1$ and $\gcd(b^m+1,b^m-1)=1$, we have $\gcd(P,p_1p_2\dotsb p_t)=1$. Next, consider $q_1,q_2,\dotsc,q_\tau$ be all prime factors of $b^{m-1}+1$. Hence, $b^{m-1}\equiv-1\pmod{q_1q_2\dotsb q_\tau}$. Let $r_1,r_2,\dotsc,r_{(P-b+1)/2}$ be positive integers, where $r_{i+1}-r_i\geq m+1$ for all $1\leq i\leq(P-b-1)/2$, be defined as follows.
\begin{itemize}
\item If $(P-b+1)/2$ is odd, then
$$b^{r_i+2}\equiv\begin{cases}
-1\pmod{q_1q_2\dotsb q_\tau}&\text{if }1\leq i\leq(P-b-1)/4;\\
1\pmod{q_1q_2\dotsb q_\tau}&\text{otherwise}.
\end{cases}$$
\item If $(P-b+1)/2$ is even, then
$$b^{r_i+2}\equiv\begin{cases}
-1\pmod{q_1q_2\dotsb q_\tau}&\text{if }1\leq i\leq(P-b-3)/4;\\
1\pmod{q_1q_2\dotsb q_\tau}&\text{if }(P-b+1)/4\leq i\leq(P-b-3)/2;\\
b\pmod{q_1q_2\dotsb q_\tau}&\text{otherwise}.
\end{cases}$$
\end{itemize}
Now, let $c=\sum_{i=1}^{(P-b+1)/2}b^{r_i}(b^m+1)$. Since $b\mid c$, we have $b\nmid(c+1)$. Next, $s_b(c)=P-b+1$ from our construction, thus $s_b(c)+b-1=P=b^m+1$, which is a factor of $c$. Recalling that $P$ is relatively prime to all positive integers up to $2b$, we have $\gcd(cb^2+j(b-1),P)=1$ for all $1\leq j\leq 2b$. It remains to prove that $\gcd(cb^2+j(b-1),s_b(c)+2(b-1))=1$ for $j\in\{b+1,2b+1\}$. Note that $s_b(c)+2(b-1)=P-b+1+2(b-1)=P+b-1=b(b^{m-1}+1)$. For any prime factor $q$ of $b^{m-1}+1$, we clearly have $q\nmid b$. Moreover, $q\nmid(b-1)$, or otherwise, $q\mid b(b^{m-1}+1)$ and $q\mid(b-1)$ imply $q\mid P$, contradicting that $P$ is relatively prime to all positive integers up to $2b$. If $(P-b+1)/2$ is odd, then
\begin{align*}
    cb^2+(b+1)(b-1)&=\left(\sum_{i=1}^{(P-b+1)/2}b^{r_i+2}\right)P+b^2-1\\
    &\equiv P+b^2-1\\
    &\equiv P+b^2-1-(P+b-1)\\
    &\equiv b(b-1)\\
    &\not\equiv0\pmod{q}
\end{align*}
and $cb^2+(2b+1)(b-1)=cb^2+(b+1)(b-1)+b(b-1)\equiv2b(b-1)\not\equiv0\pmod{q}$. If $(P-b+1)/2$ is even, then 
\begin{align*}
    cb^2+(b+1)(b-1)
    &=\left(\sum_{i=1}^{(P-b+1)/2}b^{r_i+2}\right)P+b^2-1\\
    &\equiv 2bP+b^2-1\\
    &\equiv 2bP+b^2-1-2b(P+b-1)\\
    &\equiv-(b-1)^2\\
    &\not\equiv0\pmod{q}
\end{align*}
and $cb^2+(2b+1)(b-1)=cb^2+(b+1)(b-1)+b(b-1)\equiv-(b-1)^2+b(b-1)\equiv b-1\not\equiv0\pmod{q}$. Finally, our proof is completed by noticing that $\gcd(cb^2+j(b-1),b)=1$ for $j\in\{b+1,2b+1\}$.
\end{proof}

\end{document}
